% 第九卷 12.6.2.2 节：Transformer 翻译模型（编码器—解码器）。
% 左列编码器、右列解码器，主链路自下而上；灰字为省略批维的张量形状。
% 蓝色：注意力子层；黄色：前馈子层；红色：输出；白底灰框：嵌入与位置编码。
\begin{tikzpicture}[
  x=1mm,y=1mm,
  every node/.style={font=\figurefont\footnotesize,text=FigureInk},
  op/.style={rectangle,rounded corners=1.8pt,draw=FigureStroke,line width=0.6pt,fill=white,
    minimum height=7.5mm,minimum width=36mm,inner sep=1.2mm,align=center},
  attnop/.style={op,draw=FigureBlue,fill=FigureBlueFill},
  ffnop/.style={op,draw=FigureYellow,fill=FigureYellowFill},
  outop/.style={op,draw=FigureRed,fill=FigureRedFill},
  tok/.style={font=\figurefont\footnotesize,text=FigureInk,inner sep=0.8pt},
  flow/.style={-{Stealth[length=1.8mm,width=1.3mm]},draw=FigureStroke,line width=0.7pt},
  kv/.style={-{Stealth[length=1.8mm,width=1.3mm]},draw=FigureBlue,line width=0.9pt,rounded corners=2pt},
  dim/.style={font=\figurefont\scriptsize,text=FigureMuted,inner sep=0.8pt},
  panel/.style={rounded corners=2.5pt,draw=FigureGrid,line width=0.5pt,fill=FigurePanel},
  ptitle/.style={font=\figurefont\scriptsize\bfseries,text=FigureMuted}
]
  % 层堆叠背景。
  \draw[panel] (0,20) rectangle (44,51);
  \draw[panel] (58,20) rectangle (102,70.5);
  \node[ptitle,anchor=north west] at (0.8,50.2) {编码器层$\times L$};
  \node[ptitle,anchor=north west] at (58.8,69.7) {解码器层$\times L$};

  % 编码器。
  \node[tok] (src) at (22,0) {源句：他\ 不\ 参加\ 会议};
  \node[op] (eemb) at (22,11) {$\sqrt d\,\bm E[x]+\bm P$};
  \node[attnop] (esa) at (22,28) {多头自注意力（双向）\\残差与层归一化};
  \node[ffnop] (effn) at (22,42) {前馈网络\\残差与层归一化};
  \node[op,draw=none,fill=none] (henc) at (22,60) {编码器输出$\bm H_{\mathrm{enc}}$};
  \draw[flow] (src) -- (eemb);
  \draw[flow] (eemb) -- (esa);
  \draw[flow] (esa) -- (effn);
  \draw[flow] (effn) -- (henc);
  \node[dim,anchor=west] at (23,17.5) {$S\times d$};
  \node[dim,anchor=west] at (23,54.5) {$S\times d$};
  \node[dim,anchor=east] at (21,21.8) {权重$S\times S$};

  % 解码器。
  \node[tok] (tgt) at (80,0) {右移目标：BOS\ He\ will\ not};
  \node[op] (demb) at (80,11) {$\sqrt d\,\bm E[y_{<t}]+\bm P$};
  \node[attnop] (dsa) at (80,28) {因果自注意力\\残差与层归一化};
  \node[attnop] (dca) at (80,44) {交叉注意力\\残差与层归一化};
  \node[ffnop] (dffn) at (80,59) {前馈网络\\残差与层归一化};
  \node[outop] (out) at (80,78) {乘$\bm E^{\top}$，softmax};
  \node[tok] (pred) at (80,90) {预测：He\ will\ not\ attend};
  \draw[flow] (tgt) -- (demb);
  \draw[flow] (demb) -- (dsa);
  \draw[flow] (dsa) -- node[dim,right=0.4mm,pos=0.5] {查询} (dca);
  \draw[flow] (dca) -- (dffn);
  \draw[flow] (dffn) -- (out);
  \draw[flow] (out) -- (pred);
  \node[dim,anchor=west] at (81,17.5) {$T\times d$};
  \node[dim,anchor=east] at (79,21.8) {权重$T\times T$};
  \node[dim,anchor=east] at (79,37.2) {权重$T\times S$};
  \node[dim,anchor=west] at (81,72.5) {$T\times d$};
  \node[dim,anchor=west] at (81,84) {$T\times V$};

  % 编码器输出送入每个解码层的交叉注意力。
  \draw[kv] (henc.east) -- (51,60) -- (51,44) -- (dca.west);
  \node[dim,anchor=south,text=FigureBlue,align=center] at (44.5,60.6) {键、值\\各层共用};
\end{tikzpicture}
